\section{Auditing source and rendered output}
\label{sec:latex-style-source-render}

An audit of rendered output includes verifying that every rendered glyph is
the intended character; the normative catalogue of mathematical symbol
characters is \cite{Unicode2025}.

The following small example exhibits the whole path from semantic source to a
rendered claim.  The source deliberately contains a typed signature, quantified
semantic parameters, semantic macros, and a stable label:
\begin{lstlisting}[language={[LaTeX]TeX},caption={Source for a semantic implication}]
\begin{theorem}\label{thm:sample-semantic-implication}
Let \(\Lang\) contain unary relation symbols \(P,Q\).
For every \(\Lang\)-structure \(\mathcal M\) and assignment \(s\),
if \(\mathcal M,s\entails P(x)\) and
\(\mathcal M,s\entails\forall y(P(y)\to Q(y))\), then
\(\mathcal M,s\entails Q(x)\).
\end{theorem}
\begin{proof}
By the universal clause, if \(P^{\mathcal M}(s(x))\) then
\(Q^{\mathcal M}(s(x))\).  The first hypothesis gives
\(P^{\mathcal M}(s(x))\), so \(Q^{\mathcal M}(s(x))\), which is exactly
\(\mathcal M,s\entails Q(x)\).
\end{proof}
\end{lstlisting}

The rendered statement must visibly read as a claim about every
\(\Lang\)-structure and assignment; both hypotheses must remain inside the
conditional, and \(\entails\) must render as semantic satisfaction rather than
derivability.  The rendered proof must expose the instance \(s(x)\), because
that calculation is the bridge between the universal semantic clause and the
conclusion.

Audit this example in the following order.
\begin{enumerate}
  \item \emph{Source audit.}  Confirm that \(P,Q\) are relation symbols,
        \(\mathcal M\) is a structure, \(s\) is an assignment, and \(x,y\) are
        object-language variables.  Confirm that the label is unique and that
        a later reference uses \verb|\cref| rather than a typed theorem number.
  \item \emph{Formal-logical audit.}  Expand the satisfaction clauses.  The
        hypotheses become \(P^{\mathcal M}(s(x))\) and
        \(\forall a\in\abs{\mathcal M}
        (P^{\mathcal M}(a)\to Q^{\mathcal M}(a))\); instantiation at \(s(x)\)
        yields the conclusion.  This independently checks the prose proof.
  \item \emph{Compilation audit.}  Compile enough times to resolve the label
        and bibliography.  Reject undefined references, missing citations,
        duplicate labels, malformed environments, and any overfull box that
        obscures a symbol.
  \item \emph{Rendered-page audit.}  Inspect the theorem number, line breaks,
        subscripts, turnstile/model glyphs, proof-ending mark, and surrounding
        page break.  Verify that no line makes the quantifier appear to govern
        less text than the source says it governs.
  \item \emph{Extraction audit.}  Extract text from the final PDF and check that
        the theorem heading, label target, hypotheses, and conclusion remain
        searchable.  Text extraction supplements visual inspection; it does
        not replace it.
\end{enumerate}
Passing only the source or only the rendered page is insufficient: the audit
requires agreement among the formal claim, the LaTeX source, compilation
records, cross-reference target, and visible output.

